\documentclass{situationist-preprint}

\title{Identity in Ordinary Cases}
\msid{PID-01-c}
\msfamily{PID-01}
\msdiscipline{Personal Identity and the Self}
\mstype{Consequence}
\mskeywords{personal identity; psychological continuity; bodily criterion; reidentification; sleep; phase sortal}

\begin{document}
\maketitle

\begin{abstract}
The debate over personal identity is conducted through puzzle cases: fission, teletransportation, uploading. One may worry that a criterion built to handle such cases leaves ordinary judgments of identity insecure. We derive from the answer of PID-01, psychological continuity with bodily continuity as a tie-breaker in branching, that in ordinary cases, in which a human organism's psychology continues in that organism without branching, the psychological, bodily and biological criteria coincide and identity holds determinately. It follows that the ordinary evidence for identity, bodily and behavioral, is reliable evidence for the relation that constitutes it, and that dreamless sleep, anesthesia, slow change of character and the passage from childhood are not threats to it.
\end{abstract}

\section{Ordinary cases}

\begin{definition}\label{d:ord}
A pair of person-stages $x$ and $y$, with $x$ earlier, forms an \emph{ordinary case} if both are stages of one living human organism, the organism has the capacities for rational and self-conscious thought at both times, and no person other than the one at $y$ is, at any time, psychologically continuous with $x$ through a cause of the right kind.
\end{definition}

Almost every judgment of identity anyone makes concerns an ordinary case. The answer of PID-01 is that a person persists by psychological continuity, non-branching and through a cause of the right kind, including connections of character and value, where the cause is of the right kind only if the later stages are realized in a thing of the same kind as the earlier ones; bodily continuity breaks ties when psychology branches; and \enquote{person} is a phase sortal of the organism.

\section{The consequence}

\begin{lemma}\label{l:cont}
In an ordinary case, $y$ is psychologically continuous with $x$ through a cause of the right kind.
\end{lemma}

\begin{proof}
Between any two times in the organism's personal phase, its psychological states are sustained by one brain in one living body, and each day's beliefs, desires, intentions, values and traits are strongly connected with the next day's, since nothing has intervened to disrupt them. Overlapping chains of such connections therefore run from $x$ to $y$. Their cause is the normal operation of one brain in one organism, which realizes $x$ and $y$ alike, so the kind condition is met.
\end{proof}

\begin{theorem}\label{t:main}
In an ordinary case, $x$ and $y$ are stages of the same person; and the psychological criterion, the biological criterion of PID-05 and the criterion of bodily continuity all give this verdict.
\end{theorem}

\begin{proof}
By Lemma~\ref{l:cont} and the non-branching clause of Definition~\ref{d:ord}, the psychological criterion of PID-01 is satisfied, and the tie-breaker is idle. Since $x$ and $y$ are stages of one living organism, the biological criterion is satisfied (PID-05-a, Assumption~2.7), and the same body persists from $x$ to $y$. Since \enquote{person} is a phase sortal of the organism, the person at $x$ and the person at $y$ are the organism in its personal phase, as PID-05-a, Proposition~4.1, states for every case in which an organism's psychology continues in the organism.
\end{proof}

\begin{corollary}\label{c:evidence}
In ordinary cases, bodily and behavioral evidence of identity is evidence of the relation that constitutes it.
\end{corollary}

\begin{proof}
We reidentify people by face, voice, gait, location and conduct, and occasionally by memory. In an ordinary case the body that presents these features is the organism whose continuing life sustains the psychological continuity that constitutes personal identity, by Theorem~\ref{t:main}. Bodily evidence is therefore evidence of what grounds the psychological relation, not of a rival relation.
\end{proof}

This answers an old complaint against psychological criteria \citep{Williams1957}: that we do not in practice reidentify persons by their memories and could not, since memory claims must be checked against bodily evidence. The criterion is a claim about what identity consists in, not about how it is known, and in ordinary cases what is known by bodily evidence is what the criterion requires.

\section{Interruptions and slow change}

Three kinds of ordinary case are sometimes thought troublesome.

\paragraph{Dreamless sleep and anesthesia.} During dreamless sleep or general anesthesia there may be no occurrent conscious states. But the connections that make up continuity are mainly dispositional: beliefs, intentions, values and traits of character persist through sleep in the brain that realizes them, and memories formed before are retrievable after. Continuity requires the persistence of the states and capacities, not their uninterrupted exercise \citep[ch.~10]{Parfit1984}. The sleeper and the waker form an ordinary case.

\paragraph{Slow change of character.} A person may come, over decades, to have almost none of the beliefs, values and traits she once had. The chains of connections overlap nonetheless, and identity holds, while connectedness between the distant stages is weak. This is the conclusion of PID-07 for the childhood self. What weakens is the relation that carries most of what matters (PID-06), not identity.

\paragraph{Loss of memory.} Since continuity includes connections of character and value and of other psychological capacities, the loss of autobiographical memory does not by itself break it, as PID-08 argues.

\begin{objection}
Puzzle cases show that the concept of personal identity has no determinate application in some situations. If it is indeterminate there, it may be indeterminate in ordinary cases too, since the ordinary cases differ from the puzzles only in degree.
\end{objection}

\begin{reply}
Ordinary cases differ from the puzzles in kind. In every puzzle case some relation that normally accompanies psychological continuity is absent: the organism, the body, uniqueness, or the same kind of realizer. In an ordinary case all are present together, and Theorem~\ref{t:main} shows that every criterion in contention agrees. Whatever indeterminacy the puzzles may reveal arises from the separation of relations that ordinary cases keep together, and so cannot spread to them.
\end{reply}

\section{Conclusion}

The answer of PID-01 is shaped by cases that never occur, but it is not hostage to them. In the cases that do occur, it delivers determinate identity, it agrees with every rival criterion, and it makes the evidence on which we actually rely evidence of what identity consists in.

\begin{thebibliography}{99}\small
\bibitem[Parfit(1984)]{Parfit1984} Parfit, D. (1984). \emph{Reasons and Persons}. Oxford University Press.
\bibitem[Williams(1957)]{Williams1957} Williams, B. (1957). Personal identity and individuation. \emph{Proceedings of the Aristotelian Society}, 57, 229--252.
\end{thebibliography}

\end{document}
